\documentclass[a4paper,12pt]{article}

% Packages
\usepackage{graphicx}
\usepackage{amsmath}
\usepackage{geometry}
\usepackage{fancyhdr}
\usepackage{setspace}
\usepackage{titlesec}  % For title formatting
\usepackage{booktabs}
\usepackage[hidelinks]{hyperref}
\usepackage{float}

\geometry{margin=1in}
\setstretch{1.5}

% Header and Footer
\pagestyle{fancy}
\fancyhf{}
\fancyhead[L]{Comparative Data Mining Group Project}
\fancyhead[R]{\thepage}

% 标题格式
\titleformat{\section}{\normalfont\Large\bfseries}{\thesection}{1em}{}

% 封面页
\title{
    \vspace{2cm} % 保持图片间距不变
    \centering
    \includegraphics[width=1\textwidth]{LU.jpg} \\ % 让图片居中，并调整大小
    \vspace{1cm} % 保持图片间距不变
    \textbf{\Huge Comparative Data Mining on German Credit and Thyroid Disease} \\
    \vspace{1cm} % 保持图片间距不变
    \large CDS535 - Data Mining \\
    \vspace{0.5cm} % 保持图片间距不变
    \large 2026
}
\author{
    \begin{minipage}{\textwidth}
        \centering
        Team name: to be confirmed \\
        Member A (student ID to be confirmed) \\
        Member B (student ID to be confirmed) \\
        Member C (student ID to be confirmed) \\
        Member D (student ID to be confirmed) \\
        Member E (student ID to be confirmed)
    \end{minipage} \\[25pt]
    \begin{minipage}{\textwidth}
        \centering
        \textbf{GitHub:} \url{https://github.com/chenyanwu83-blip/DM-Group-Project1}
    \end{minipage}
}
\date{}




\begin{document}

% Title Page
\maketitle
\thispagestyle{empty}
\newpage
\thispagestyle{empty}
\begin{abstract}
% TODO(E): Replace with a 150--200 word abstract after D supplies verified results.
\noindent This report compares data mining workflows on the UCI Statlog German Credit Data and the sick subset of the UCI Thyroid Disease collection. The final abstract will summarize the methods and verified findings after the experiments are complete.
\newline

\noindent \textbf{Keywords}: classification, class imbalance, association rules, UCI datasets
\end{abstract}
\newpage
\thispagestyle{empty}
\tableofcontents
\thispagestyle{empty}
\newpage





% Table of Contents
% Start page numbering from the Table of Contents
\setcounter{page}{1}  % Start counting from 1

\section{Group Project Requirement}
\label{sec:requirement}

This group studies two datasets from the UCI Machine Learning Repository: Statlog German Credit Data and the sick subset of the Thyroid Disease collection~\cite{dua2019uci,uci_credit_g_1994,uci_thyroid_1986}. Each dataset will follow a documented preprocessing, modeling, and evaluation workflow. The report will compare the datasets and include four supervised algorithms, an association-rule component, cross-validation, an independent evaluation protocol, figures, and a reproducible comparison table.

% TODO(E): Add the confirmed five-member task division and contribution table.
% Check the final report against every row of Table 3 in the teacher's requirements.

\section{Introduction}
\label{sec:introduction}

\subsection{Background and Motivation}
Automated credit scoring supports lending decisions in which the two error types are
economically asymmetric: approving an applicant who turns out to default (a missed
\texttt{bad} case) destroys principal, whereas rejecting a creditworthy applicant (a
false \texttt{bad} alarm) merely forgoes interest income. The Statlog German Credit
data makes this asymmetry explicit through an official 5:1 misclassification-cost
matrix, in which missing a \texttt{bad} applicant costs five times as much as
rejecting a \texttt{good} one~\cite{uci_credit_g_1994}. Class imbalance amplifies the
problem: \texttt{bad} applicants form only 30\% of \texttt{credit-g} and about 6\% of
the \texttt{sick} thyroid subset~\cite{uci_thyroid_1986}, so accuracy is dominated by
the majority class --- a trivial rule that always predicts \texttt{good} already
achieves 0.70 accuracy on \texttt{credit-g} while never identifying a single default.
Practitioners can draw on three families of remedies --- synthetic minority
over-sampling~\cite{chawla2002smote}, cost-sensitive training, and decision-threshold
adjustment --- yet their relative merits change with the model family and the degree
of imbalance, and are rarely compared under one leakage-safe protocol with an
application-grounded cost measure.

\subsection{Research Significance}
This project delivers such a comparison on two UCI datasets~\cite{dua2019uci} that
differ deliberately in scale, imbalance ratio, missingness and domain. For risk
management, it shows which imbalance intervention actually minimises the official
misclassification cost rather than a surrogate metric. For methodology, it
demonstrates a reproducible experimental discipline: preprocessing specifications
frozen before the main matrix, locked data splits, SHA-256 provenance chained from
raw data to every result file, and all imputation, encoding, scaling and resampling
fitted strictly inside the cross-validation folds, so that any observed gain is
attributable to the strategies rather than to information leakage. For
interpretability, association rules mined from the same data are triangulated
against tree-based model explanations. The full pipeline is implemented with
scikit-learn~\cite{pedregosa2011sklearn}, imbalanced-learn~\cite{lemaitre2017imblearn}
and mlxtend~\cite{raschka2018mlxtend}, with Weka~\cite{frank2016weka} as the course
reference tool.

\subsection{Research Questions and Contributions}
The study is organised around four research questions:
\begin{itemize}
  \item \textbf{RQ1.} How does the degree of class imbalance moderate the gains of
  imbalance-handling strategies? The two datasets provide a 70:30 and a 94:6
  endpoint for this question.
  \item \textbf{RQ2.} Under what conditions does accuracy mislead model selection on
  imbalanced data, and which metrics (positive-class F1, ROC-AUC, official
  misclassification cost) track decision utility instead?
  \item \textbf{RQ3.} Under the official 5:1 cost matrix, does the best
  model--strategy pairing change relative to the rankings obtained by accuracy or F1?
  \item \textbf{RQ4.} How does the encoding of nominal attributes (one-hot vs.\
  ordinal) interact with the model family (linear/kernel vs.\ tree-based) on
  mixed-attribute data?
\end{itemize}

\noindent The main contributions are: (i) a paired two-dataset design contrasting
five dimensions at once --- scale (1{,}000 vs.\ 3{,}772 records), imbalance (70:30
vs.\ 94:6), missingness (0 vs.\ about 16\%), attribute composition and domain ---
which directly probes RQ1; (ii) a full experiment matrix of four classifiers
(decision tree, random forest, logistic regression, RBF-SVM) $\times$ four
strategies (plain, SMOTE, cost-sensitive, threshold adjustment) built on
leakage-safe pipelines with locked splits; (iii) the official cost matrix treated
as a first-class evaluation axis alongside F1 and ROC-AUC (RQ2, RQ3); (iv) a
controlled encoding ablation with a supplementary attribution control that
separates encoding from scaling effects (RQ4); and (v) an interpretability
triangulation of Apriori rules, decision-tree paths and random-forest feature
importances (Section~\ref{sec:apriori}). Statistical comparisons use fold-level
tests with the Nadeau--Bengio
correction~\cite{nadeau2003inference,demsar2006comparison}. Experimental results
are reported in Section~\ref{sec:results}.

\section{Data Distribution and Preprocessing}
\label{sec:data}
\subsection{Statlog German Credit Data}
% A owns this subsection: UCI provenance, schema, class distribution, and preprocessing.

\subsection{Sick Dataset}
% B owns this subsection: UCI provenance, label meaning, missingness, and preprocessing.

\section{Basic Concepts of Algorithms}
% C owns this file: classifiers and imbalance strategies.

\section{Model Training}
% C owns this file: pipelines, fixed settings, cross-validation, and holdout protocol.

\section{Results and Evaluation}
\label{sec:results}
% D owns this file: metrics, figures, cross-dataset comparison, and statistical analysis.

\subsection{Association Rule Analysis}
\label{sec:apriori}

% E: Insert the verified credit-g Apriori data preparation and parameter settings.
% Report representative class=bad rules with support, confidence, and lift.
% Interpret the rules as associations, not causal effects.
% Compare selected patterns with C's decision tree and random forest outputs.

\section{Conclusion and Future Work}
% D drafts this file; E helps align conclusions with measured results and limitations.


% Bibliography (if required)
\bibliographystyle{plain}
\bibliography{references}

\clearpage
\section*{Team Contributions}
\label{sec:contributions}

% E: Replace this comment with the teacher's Table 5 fields:
% member name and student ID, role, specific contribution, workload percentage.
% Confirm every entry and the total workload allocation with the group.


\end{document}
